% Insert into report_las_sto3g.tex immediately AFTER the subsection
% "Cost functions and the reference state" (i.e. after Eq.~\eqref{eq:commutator}
% and its paragraph), and BEFORE "Candidate libraries and the ranked scan".
%
% It promotes the claim already made informally at the end of the subsection
% "Parity rows, exact group, quotient" -- that c(g+e)=c(g) -- to a stated and
% proved proposition, and then draws the consequence that governs Q3 and Q4.

\subsection{The cost is a class function on the quotient}
\label{sec:classfunction}

The exclusion of exact-dressed duplicates in Eq.~\eqref{eq:ranktest} was
justified above by the assertion $c(g+e)=c(g)$. That identity is not a
convenience: it determines what any cost-ranked selection rule is able to
choose, and it is therefore stated and proved here.

\begin{proposition}[Coset invariance of the cost]
\label{prop:coset}
Let $|\Psi\rangle$ lie in a fixed exact sector, so that $Z(e)|\Psi\rangle
=\sigma_e|\Psi\rangle$ with $\sigma_e\in\{+1,-1\}$ for every $e\in\Ecal$. Then
for all $g\in\FF^{N}$ and all $e\in\Ecal$,
\begin{equation}
c^{\mathrm{NC}}_{\Psi,U}(g+e)=c^{\mathrm{NC}}_{\Psi,U}(g),
\qquad
c^{\mathrm{var}}_{\Psi,U}(g+e)=c^{\mathrm{var}}_{\Psi,U}(g).
\label{eq:cosetinvariance}
\end{equation}
\end{proposition}

\begin{proof}
Addition of rows over $\FF$ is multiplication of the operators in
Eq.~\eqref{eq:Zg}, so $Z(g+e)=Z(g)Z(e)$. By Eq.~\eqref{eq:exact},
$[\Ht(U),Z(e)]=0$, hence
\begin{equation}
[\Ht(U),Z(g)Z(e)]
=[\Ht(U),Z(g)]\,Z(e)+Z(g)\,[\Ht(U),Z(e)]
=[\Ht(U),Z(g)]\,Z(e).
\label{eq:commexpand}
\end{equation}
Applying this to $|\Psi\rangle$ and using $Z(e)|\Psi\rangle=\sigma_e|\Psi\rangle$
gives $[\Ht,Z(g+e)]|\Psi\rangle=\sigma_e\,[\Ht,Z(g)]|\Psi\rangle$. Since
$\sigma_e^{2}=1$, the norms in Eq.~\eqref{eq:nc} coincide. For
Eq.~\eqref{eq:var}, $\langle\Psi|Z(g)Z(e)|\Psi\rangle
=\sigma_e\langle\Psi|Z(g)|\Psi\rangle$, and the expectation is squared.
\end{proof}

Both cost functions are therefore constant on the cosets $g+\Ecal$: they are
well defined on the quotient $\Qcal$ of Eq.~\eqref{eq:quotient}, not merely on
$\FF^{N}$. Write $\bar c(\bar g)$ for the induced function, $\bar g$ denoting
the class of $g$.

Two structural consequences follow, and together they account for the
selection-rule results of Sections~\ref{sec:q3} and~\ref{sec:q4}.

\paragraph{(0) The ground set is the spatial subregister, not $\FF^{N}$.}
\begin{sloppypar}
Candidates are generated as products of the spin-symmetric pairs
$\Zb_p=Z_{p\alpha}Z_{p\beta}$, i.e.\ as images of the embedding $\iota$ of
Eq.~\eqref{eq:iota}, so the ground set of the matroid below is
$\iota(\FF^{n})\subset\FF^{N}$ and not $\FF^{N}$ itself. This is an
implementation restriction, not a requirement of the procedure: the candidate
generator, the GF(2) rank test and the cost evaluation all act on rows of
length $n$, and spin-resolved rows such as $Z_{p\alpha}$ alone are never
formed. Its consequences are that the accessible budget is bounded by
$n-\rsp$ rather than $N-r$ (Eq.~\eqref{eq:budget-sp}), and that the optima
reported here are optima over $\iota(\FF^{n})$ only; a spin-resolved search
could in principle reach a lower selection cost at the same $M$. Removing the
restriction requires the candidate library, the rank test and the cost to be
carried in the length-$N$ register throughout --- a change of representation
rather than of algorithm --- and it has not been done here. All statements in
Sections~\ref{sec:q3}, \ref{sec:q4} and~\ref{sec:budget} are therefore to be
read as holding within the spin-symmetric ground set.
\end{sloppypar}

\paragraph{(i) The pool is an independent set of a linear matroid.}
The acceptance test of Eq.~\eqref{eq:ranktest} admits $g$ exactly when
$\bar g$ is linearly independent of the classes already accepted. The
admissible pools of size $M$ are therefore precisely the independent sets of
size $M$ of the linear matroid on $\Qcal$, and the budget bound
Eq.~\eqref{eq:budget} is the statement that the rank of that matroid is
$\dim\Qcal$.

\paragraph{(ii) Greedy ranking is optimal at every budget.}
For a matroid with non-negative element weights, the greedy algorithm returns a
minimum-weight independent set of size $M$ for \emph{every} $M$, not only at
full rank. Since the weight $\bar c$ is well defined on $\Qcal$ by
Proposition~\ref{prop:coset}, any two selection rules that both attain the
greedy optimum select independent sets of equal total weight, and hence
subspaces $\operatorname{span}\{\bar g_1,\dots,\bar g_M\}\subseteq\Qcal$ that
are optimal in the same sense. Minimum-weight independent sets need not be
unique; when they are not, the rules may return different rows while spanning
the same subspace.

\paragraph{Gauge freedom.}
A rule reports rows $g_j\in\FF^{N}$, but only the classes $\bar g_j$ carry
information: by Proposition~\ref{prop:coset} the cost is unchanged, and by
Eq.~\eqref{eq:ranktest} the accepted rank is unchanged, if any $g_j$ is
replaced by $g_j+e$. The group generated by $\Ecal$ together with the pool,
and therefore the sector partition, the block dimensions $\Dmax$ and the
coupled dimension $K$, depend only on
$\operatorname{span}\{\bar g_1,\dots,\bar g_M\}$. Differences between rules in
the reported rows are consequently \emph{not} evidence of a different
selection; only a difference in the spanned subspace is. This distinction is
applied throughout Section~\ref{sec:q4}, where rules are compared by span
rather than by row equality.

\paragraph{A degenerate budget.}
If $M=\dim\Qcal$ then any $M$ independent classes form a basis of $\Qcal$, so
$\operatorname{span}\{\bar g_1,\dots,\bar g_M\}=\Qcal$ for \emph{every}
admissible pool, and
$\operatorname{span}(\Ecal\cup\Gcal)=\FF^{N}$ irrespective of the rule. At that
budget the partition is rule-independent by construction and carries no
discriminating information; a rule can influence $K$ only indirectly, through
the orbital objective of Eq.~\eqref{eq:objective}. This degeneracy is the one
noted in Section~\ref{sec:matroid} for the span test; the statement here is
its general form, and it applies to $\Dmax$ and the block count as well.
Comparisons at $M=\dim\Qcal$ must therefore be read with the degeneracy in
mind, and Section~\ref{sec:budget} repeats the campaign at $M=\dim\Qcal-1$,
where the spanned subspace is proper and the degeneracy is absent.
